\section{Discussion}
\label{sec:discussion}

\paragraph{Rewards and residuals.} A residual-based loss and a reward are different objects. The reward is part of the economic model; a residual loss is a consistency criterion that a researcher derives from the model. The DEQN of the companion thesis illustrates the risk: with an envelope condition that omitted the direct effect of capital on utility and first-order conditions that counted the effect on future value twice, it converged, with a vanishing loss, to a consumption rate of 0.53--0.93 instead of 0.74. The actor-critic has no such failure mode. Its objective is built from the reward and the laws of motion, and the conditions that correspond to the envelope conditions, the gradient targets of the critic, are generated by automatic differentiation of the Bellman equation.

\paragraph{The SCC is a small gradient of a large function.} The ablations show that the obstacle for a reward-based method is not the policy but the shadow price. A critic trained only on values can be very accurate and still imply the wrong sign for the SCC, because the marginal value of carbon is four orders of magnitude smaller than the value itself and is confounded with time along the simulated trajectory. Both DEQNs and the actor-critic must therefore give the carbon conditions extra weight. In the DEQN, this requires hand-chosen scales for the costate outputs and residuals. In the actor-critic, it follows from normalizing each component of the gradient loss by the running scale of its own target, which requires no knowledge of the economics. Value-gradient learning is, in this sense, the natural critic for economic problems, in which policies are driven by shadow prices.

\paragraph{When residuals are still preferable.} The correctly specified DEQN remains about an order of magnitude more precise on this problem, and residual-based methods have a decisive advantage beyond precision: they apply to equilibria that do not solve any single optimization problem. In decentralized economies with market clearing, heterogeneous agents or externalities that individual agents ignore, there is no single reward whose maximum is the equilibrium, but the equilibrium conditions can still be written as residuals. The reward-based approach is therefore a complement for problems with a planner, or with agents whose optimization problems can be learned separately, rather than a replacement.

\paragraph{Fragility.} The method has failure modes of its own. The actor reacts to the critic's gradient within tens of iterations, while the critic's carbon gradient takes hundreds of iterations to converge. If the actor follows a wrongly signed critic long enough to saturate its output, it does not recover, even after the critic becomes correct (Section \ref{sec:ablation}). The critic warm-up and the delayed actor updates address this race, but they are hyperparameters. I chose them in trial runs on the DICE 2016 model and applied them unchanged to the DICE 2023 model, and Section \ref{sec:seeds} reports how robust the outcome is to the random seed. Saturation could be avoided altogether by a parameterization of the actor that does not saturate, or by penalizing extreme outputs, which I leave for future work.

\paragraph{Limitations and extensions.} The model is deterministic, the dynamics are known and differentiable, and business-as-usual emissions follow an exogenous path so that the benchmark is exact. Each of these assumptions suggests an extension. First, with known differentiable dynamics, a critic is not strictly necessary: the policy can be trained by differentiating through simulated trajectories, with a learned value only at the end of a truncated horizon. Second, for stochastic or non-differentiable climate models, for instance with tipping points or with black-box emulators, model-free methods such as PPO or GRPO would become the relevant comparison. They estimate the policy gradient from sampled returns, and the small magnitude of the carbon signal suggests that they would need very large samples, which a comparison on this benchmark could quantify. Third, with endogenous business-as-usual emissions, the problem no longer has a closed form, but the actor-critic applies unchanged.
